\section{Discussion}
\label{sec:discussion}

\subsection{What the framework resolves}

The comparisons in Section~\ref{sec:tests} indicate that a one-dimensional
formulation that conserves gas mass and momentum in the shaft and pipe, and
that treats trapped pockets by an equation of state on the resolved pocket
volume, is sufficient to select the geysering regime without regime-specific
inputs. In the wide-tower test the venting pathway---pocket discharge
through the tower with counter-current exchange at the mouth---sets a
pressure plateau whose level and duration are controlled by the pocket
volume and the tower hydrostatic column; both are reproduced because they
follow from conservation rather than from calibrated closures. The interface
climb velocity in the wide tower is likewise reproduced at the
Taylor-bubble scale without any drift-flux input, because the buoyant rise
of the gas core through the resolved liquid column is part of the two-fluid
solution.

The regime discrimination between Tests A and B is controlled by the tower
diameter through two competing rates that the model resolves: the rate at
which the pocket discharges gas into the tower base, and the rate at which
the tower can pass gas through its water column and vent it at the rim. In
the wide tower the column passes the gas with a modest surface rise; in the
narrow tower the same discharged volume displaces the column to the rim.
This is the mechanism identified experimentally in
\cite{vasconcelos2011}, and it emerges here from the computed state.

Campaign 2 tests whether this discrimination generalizes beyond a single
apparatus, and the answer is structured rather than uniform. Across the
63-configuration sweep of the \cite{cong2017} parameter ranges, computed
with the same frozen solver and no knowledge of the experimental
criterion, the model reproduces both end members of the regime boundary
exactly---every wide-riser configuration ($D_r/D>0.62$) and every
small-volume configuration ($V^{*}_{air}<3.42$) vents without eruption,
and the narrowest risers with large pockets erupt---and it reproduces the
interaction sense of the criterion, since both thresholds derive from the
same resolved competition: the pocket must sustain its overpressure
against the shrinking column (volume condition) faster than the riser can
vent the gas through the column (diameter condition). Inside the
transition band, however, the model is conservative: 24 of 63
configurations that the criterion marks as geysering are computed as
strong but sub-rim column lifts ($1.0$--$1.6$~m against the $1.8$-m rim),
and the corresponding two intermediate Series-B runs (B-H2, B-H3) are the
only classification misses among the seven high-speed-camera runs. All
disagreements are one-sided---no configuration erupts in the model that
the criterion excludes---so the frozen configuration errs toward missing
marginal geysers, not toward false alarms.

\subsection{Identified limitations}

The completed campaigns expose five limitations, all at the level of the
junction and closure treatment rather than of the two-fluid formulation, and
all consistent with the model simplifications listed in the companion paper
\cite{xue2026model}.

\emph{Junction base datum.} The two branches of Campaign~1 currently use
different elevation datums for the shaft-base coupling: the wide-tower run
transmits the invert piezometric head to the shaft base, while the
narrow-tower run transmits the crown pressure (the tower physically taps
the pipe crown, one bore above the invert). The crown datum is the
geometrically correct choice, and in the narrow tower the difference,
up to $\rho_l g D=0.15L$ of spurious driving head, decides between a column
pinned at the rim and the measured quasi-static stand. In the wide tower the
distinction is masked by the aerated column (the visible level includes
gas-holdup swell) and the invert coupling reproduces the measurements; a
single junction closure that carries the vertical pressure distribution
across the mouth---rather than one scalar datum---is the consolidation
target.

\emph{Narrow-shaft climb closure.} The measured interface climb in the
narrow tower ($V^{*}_{int}=1.43$, far above the Taylor-bubble scale) is
approached only when the shaft gas connected to the pocket is treated as a
pocket-fed core whose nose is raised by the mouth influx, with the influx
bounded by the Wallis counter-current flooding scale; with the dissipative
cavity celerity the computed nose speed is $0.71$---half the measured
value, but still twice the Taylor scale---because the mouth influx inherits
the reduced crown feed. These two elements
use literature constants ($C_w=0.9$) and geometric closure levels, not
per-case fits, but they are engaged only in the driven regime
($H_{a0}>Y_{fs,0}$ with gas at the junction); the undriven regime retains
the dispersed-bubble drag path. In Campaign~2's venting-without-eruption
runs the driven-path feed produces the climb in surge pulses (fastest
sustained climbs $0.7$--$1.2$~m/s) where the measurements show a steady
Taylor-scale ascent ($0.42$--$0.48$~m/s); a vertical-annular entrainment
closure conditioned on the local gas flux would unify the two paths and
smooth the pulsing.

\emph{Crown-cavity celerity and the arrival/venting trade-off.} Two
elements set the timing of the gas-phase sequence. First, a discrete bias
in the front-tracking hand-off was identified during this work and removed:
releasing the cavity nose cell at $95\%$ of the region-mean layer depth let
the detected front advance on $95\%$ of its volume budget, biasing the
transit to $1.09$ times the nominal celerity. Second, the cavity celerity
coefficient itself: Benjamin's $0.542\sqrt{gD}$ is the energy-conserving
upper bound, while real (dissipative) cavities run below it, and the
middle-pipe transits of the three experimental repetitions correspond to
$0.48$--$0.52\,\sqrt{gD}$. The frozen configuration uses the dissipative
value $0.48$ (one value for both tests), which places the arrival events
inside the repetition scatter (Test~A liftoff $7.5$ against onsets
$7.3$--$7.9$; Test~B entry $-0.13$, eruption $-0.20$). The trade-off is
stated openly: the same crown current that carries the cavity nose also
feeds the vent and the narrow-tower pocket-fed core, so the reduced
celerity delays the venting-stage events (head-decay onset and interface
breakthrough late by $0.4$--$0.6$ units in both tests) and halves the
computed narrow-tower climb rate ($V^{*}_{int}=0.71$ against $1.43$). A
closure that separates the nose celerity from the vent feed rate---the
former set by the front dynamics, the latter by the pocket overpressure---%
is the identified refinement.

\emph{Film-flow pocket topology in the venting branch.} In the wide-tower
experiment the gas ascends the shaft as a single Taylor pocket that remains
pneumatically continuous with the pipe pocket: its wall film is
shear-supported, so the pocket pressure follows the weight of the water
column standing \emph{above the pocket nose}, the transducer head decays
gradually \emph{during} the climb, and the free surface creeps upward by
pocket expansion (Figs.~5 and~10 of \cite{vasconcelos2011}; their TPA model
encodes this topology by construction, as a lumped pocket with a
Batchelor-type film-flow closure). In the present field formulation the
tower gas transits as a dispersed phase interleaved with liquid, the tower
base carries the weight of the full standing column, and the computed head
therefore holds its plateau until breakthrough
(Fig.~\ref{fig:caseA_tpa}); the climb-phase free-surface creep is likewise
absent ($V^{*}_{fs}\approx0$ against the measured $0.048$). We tested a
connected-pocket closure that relaxes the pipe-pocket pressure toward the
hydrostatic environment of the tower gas nose while crediting the
transferred volume to the tower (retained as a solver option). It
reproduces the free-surface creep quantitatively ($Y^{*}_{fs}$ from $0.59$
to $0.64$ during the climb, against $0.55$--$0.65$ measured) without
degrading the validated event times, but the head history develops a
spurious hump near the catch instead of the measured monotone decay, and
more aggressive relaxation destabilizes the junction dynamics. We conclude
that the during-climb head decay is a property of the coherent-pocket
topology that a dispersed one-dimensional two-fluid closure cannot
represent stably at this grid scale---consistent with
\cite{vasconcelos2011}'s own choice of a lumped-pocket formulation for
exactly this regime---and the frozen configuration retains the
plateau-then-collapse behaviour, which preserves the plateau level, the
collapse onset time, and the interface kinematics.

\emph{Release oscillation damping.} The underdamped release-piston
oscillation visible in both tests is a genuine mode of the system---a water
column on a gas spring---and is annotated as such in the experimental
schematic of \cite{vasconcelos2011}; the experimental traces damp within
about one time unit, the model within several. The damping in the
experiment is provided by three-dimensional mixing losses at the valve and
coupling, which the constant valve and junction loss coefficients capture
in the mean but not cycle-resolved. In the narrow tower the surviving
oscillation also lets the column splash briefly to the rim during the
release swing before it settles at the quasi-static stand; in Campaign~2
the same simplification sharpens the B-H1 eruption jet ($1.71$ against
$0.92$~m/s as event speeds).

\emph{Volume-neutral mouth exchange in the reservoir-fed layout.} The
inverted-glug exchange at the tee mouth---gas up, an equal water volume
down---is what keeps the pocket pressure intact while gas transits a
standing column, and it is essential to the B-H1 geyser: without it the
narrow-riser pocket cannot rebuild the overpressure whose release ejects
the column. But in the reservoir-fed layout of Campaign~2 the same
exchange cancels part of the one-to-one column lift that the entering gas
volume should produce, and the cancellation grows with riser width: the
computed B-H6 free-surface rise is $+0.12$~m against the measured
$+0.63$~m, and across the 63-configuration sweep the under-lift accounts
for the one-sided transition-band misses (Section~\ref{sec:tests:cong2017:map}).
Gating the exchange off for reservoir-backed cases was tested and
rejected: it restores the wide-riser lift but opens the narrow pocket's
volume during the balanced venting phase, its pressure never rebuilds, and
the B-H1 geyser is lost. The two branches place conflicting demands on a
single mouth closure; a formulation that meters the exchange by the local
counter-current capacity of the mouth (rather than volume neutrality) is
the identified refinement, together with the entrainment closure above.

\subsection{Scope of the present validation}

The completed campaigns cover three apparatuses and three kinds of
evidence: time-resolved two-regime histories in a single-shaft geometry
(Campaign~1), regime classification over a systematic parameter sweep in
a riser--constant-head geometry (Campaign~2), and a prototype-motivated
junction-chamber geometry with a downstream-boundary branch pair and a
transient mass-oscillation geyser (Campaign~3), all under one frozen solver
configuration and a fully synchronous junction coupling. Within this scope
the conclusions are: in Campaign~1, regime selection is exact and both
branches are quantitative in their pressure and interface histories, with
absolute event times uniformly early by $0.5$--$0.7$ time units; in
Campaign~2, the regime boundary is reproduced in structure (both end
members, the interaction sense, the branch differentiation of the pipe-end
pressure) with a one-sided conservative bias in the transition band (5/7
of Series~B, 39/63 of the sweep, no false positives); in Campaign~3,
the A2/B3 pair flips branch with downstream boundary alone, C9 phase~1
matches the first transient geyser in timing and first-peak amplitude
against both measurement and Eq.~(\ref{eq:liu_eq7}), while compressive
spikes and pocket-driven late geysers remain limited by the frozen
elastic celerity and the reduced one-dimensional chamber treatment.
Three caveats bound the claims. First, the narrow-shaft closure elements of Test~B
(crown-datum base coupling, pocket-fed core, flooding-limited entry) are
engaged only in the driven regime; their consolidation with the wide-tower
configuration into a single junction closure is pending. Second, the
volume-neutral mouth exchange places conflicting demands on the two
Campaign-2 branches (see above), and its resolution will move the
transition-band classifications. Third, Campaign~3 uses a composite
one-dimensional chamber representation and does not resolve pocket
transport to the junction or chamber gas venting; comparison at larger
scale \cite{muller2017} remains a separate target.

\section{Conclusions}
\label{sec:conclusions}

This paper applied the decoupled two-fluid cut-cell finite-volume framework
of \cite{xue2026model} to laboratory geysering experiments, representing the
vertical shaft as the same one-dimensional two-fluid system rotated to the
vertical and coupled to the horizontal conduit through junction-loss
closures. Three campaigns were completed, all computed from the experimental
initial conditions without prescribed release rates or per-case fitting.

Against the ventilation-shaft experiments of Vasconcelos and Wright
\cite{vasconcelos2011}, the model selected the correct regime in both the
non-geysering and geysering configurations. In the non-geysering branch it
reproduced the pressure plateau level and duration, the onset of the
terminal venting decay, the maximum free-surface rise, and the air--water
interface climb velocity. In the geysering branch it reproduced the
pre-eruption quasi-static stand of the tower column ($0.84L$ against
$0.82L$), the crown-datum pressure plateau ($0.76$ against $0.76$), the
eruption to the shaft rim observed in every repetition, the interface and
free-surface climb velocities ($V^{*}_{int}=1.19$ against $1.43$;
$V^{*}_{fs}=0.42$ against $0.44$), and the sharp plateau--collapse of the
head at pocket venting; the narrow-shaft closure combines a crown-datum
base coupling, a pocket-fed gas core, and a Wallis flooding limit on the
entry flux, using geometric and literature-scale constants. The remaining
biases are a uniform early shift of the gas-phase event times
($0.3$--$0.6$ time units, within the uncertainty of the hand-operated
valve opening) and an underdamped release oscillation.

Against the horizontal-pipe/vertical-riser experiments of Cong, Chan and
Lee \cite{cong2017}, computed with the same frozen solver under the fully
synchronous coupling, the model reproduced the branch structure of the
published two-parameter geysering criterion: the geysering end member
(B-H1, with its compression-surge drive to $1.71\,H_0$), the entire
no-geysering half-plane $D_r/D>0.62$ and small-volume region, and the
branch differentiation of the pipe-end pressure that only a coupled
riser--pipe computation can produce. Its errors are one-sided and
conservative: two intermediate Series-B runs and 24 of 63 sweep
configurations that geyser in the experiment are computed as strong but
sub-rim column lifts, with no false positives; the bias is traced to the
volume-neutral mouth exchange, whose conflicting demands in the two
branches are documented as the primary closure target.

Against the junction-chamber experiments of Liu, Shao and Zhu
\cite{liu2020}, computed with the same frozen solver in a composite
one-dimensional representation, the model reproduced the downstream-boundary
branch flip between Cases A2 and B3, placed the B3 eruption on the
experimental $h$--$P_{\max}$ regression line (with acoustically
under-resolved peaks at $a_{wh}=40$~m/s), and matched the first transient
geyser of Case C9 in peak timing and magnitude against both measurement
and the authors' Eq.~(\ref{eq:liu_eq7}). Pocket-driven geysers after
pocket arrival at the chamber were declared out of scope.

These results support the use of gas-conserving one-dimensional two-fluid
models for system-scale geysering screening, where the primary question is
regime selection and venting behaviour rather than jet detail, and they
identify the junction and entrainment closures as the targets for
quantitative improvement in the geysering branch.
